\section{Conclusion}\label{sec:conclusion}

On the topic of developing a domain expert with a large language model like ChatGPT, this thesis built a hybrid multimodal lab assistant for peripheral blood smears, in which deep-learning models classify the cells in the image and a domain-expert LLM, grounded in hematology textbooks through RAG, interprets the findings. The domain expert was built in two complementary forms. The first is a GPT-4o agent that answers from a retrieval-grounded hematology knowledge base of 1{,}049 textbook passages, uses domain tools for reference ranges, cell counts and uncertainty, and returns a structured interpretation with a complete tool trace. The second is an open-weight Llama-3.1-8B model, adapted to the domain with QLoRA in under an hour on a free GPU, which can run on local hardware without a commercial API. Two image-analysis stages supply the case evidence: a YOLOv8s detector (mAP@0.50 = 0.985) and an EfficientNet-B0 classifier (98.60\,\% accuracy on a clean held-out test set), whose predictions carry Monte Carlo Dropout uncertainty --- which identifies misclassified cells with an AUROC of 0.97 --- and Grad-CAM saliency.

The central finding concerns what each technique contributes to a domain expert. Fine-tuning changed \emph{how} the model answers: its answers became concise and textbook-like, matching held-out reference answers far better than those of the base model and of GPT-4o, and it stopped inventing book references. It did not change \emph{what} the model knows: an LLM judge found no gain in clinical correctness, accuracy on external exam questions was unchanged, and the model attached a source line to incorrect answers as confidently as to correct ones. A citation generated by a language model is therefore not evidence. Knowledge and provenance must come from retrieval of real, inspectable passages, and the decision to involve a human must be taken by explicit rules. The end-to-end case study confirmed both points: the recorded tool trace made every error traceable, but the expert still turned a two-cell sample into a ``highly abnormal'' differential, because at the time no rule checked the number of cells --- a rule that has since been added.

The practical answer to the thesis topic is thus a layered design: \emph{fine-tuning for behaviour, retrieval for knowledge, tools for case data, and deterministic rules for safety}. The system is a research prototype. It has not been validated clinically, its vision models have not been tested on data from other laboratories, and its expert answers were evaluated with automatic measures rather than by hematologists.

\subsection{What was not achieved}

Several aims were met only partly, and they are stated here explicitly.

\begin{itemize}
  \item \textbf{The fine-tuned open model did not become a better expert in substance.} The hope that QLoRA fine-tuning would make Llama-3.1-8B more knowledgeable was not confirmed: clinical correctness (LLM judge) and exam accuracy were unchanged, and the model remained clearly less correct than GPT-4o. It is therefore not yet a replacement for the hosted model, only a self-hostable starting point.
  \item \textbf{The fine-tuned model was not evaluated inside the full assistant.} It was tested stand-alone, without retrieved passages; the combination of the open model with RAG and tools --- the configuration that the results suggest --- was not run end-to-end.
  \item \textbf{The end-to-end evaluation is a single case.} Only one complete run of the assistant was analysed in depth. It exposed a real failure (the two-cell differential), which led to a new safety rule, but a systematic evaluation over many smears, ideally rated by hematologists, is still missing.
  \item \textbf{The first classifier evaluation was flawed.} The original training run selected its checkpoint on the test images. This was detected during the work and corrected by retraining; the corrected accuracy (98.60\,\%) is 0.84 points lower than the originally reported figure.
  \item \textbf{Generalisation was not measured.} All image results come from curated benchmark data. A single example from another source was misclassified, which suggests that performance on smears from other laboratories may be lower.
  \item \textbf{The optional CBC input was implemented but not evaluated.}
\end{itemize}

None of these points invalidates the main conclusion; rather, they define the limits within which it holds.

\subsection{Future work}

\begin{enumerate}
  \item \textbf{Multi-field analysis.} The new minimum-cell rule means that a single field of view will almost always be flagged. Aggregating the cells of many fields (or a whole-slide image) into one differential is the natural next step towards a clinically meaningful count.
  \item \textbf{Retrieval-aware fine-tuning.} Train the open model on examples that include the retrieved passage, so that it learns to answer from and cite the evidence it is given, and evaluate it inside the agent against GPT-4o on the same cases.
  \item \textbf{Expert evaluation.} Have hematologists or medical students rate a sample of answers for correctness and faithfulness, and compare their ratings with those of the LLM judge.
  \item \textbf{Review threshold.} Decide with clinical users whether MEDIUM-uncertainty cells should also be flagged: on the test set this would catch 82\,\% instead of 51\,\% of the errors at a review rate of 4.8\,\% of cells.
  \item \textbf{External validation.} Test both vision models on smears from another source (e.g.\ Raabin-WBC or local data), ideally with a donor-level split, and repeat the calibration and error-detection analysis under this distribution shift, comparing MC-Dropout with deep ensembles.
  \item \textbf{Laboratory values.} Add a CBC input to the web interface and evaluate whether CBC values improve the interpretations.
\end{enumerate}
